\documentclass[a4paper,11pt]{ltjsarticle}
\usepackage{luatexja-fontspec}
\usepackage{amsmath,amssymb}
\usepackage{graphicx}
\usepackage{enumitem}
\usepackage{hyperref}

\title{熱伝導方程式の差分解法についての実験レポート}
\author{物理学科 3年 \quad 学籍番号 0312345}
\date{2026年10月}

\begin{document}
\maketitle

\section{目的}
一次元の熱伝導方程式を陽的差分法で解き、安定条件と誤差の関係を確かめる。

\section{理論}
温度 $u(x,t)$ は
\begin{equation}
  \frac{\partial u}{\partial t} = \kappa \frac{\partial^2 u}{\partial x^2}
  \label{eq:heat}
\end{equation}
に従う。格子幅 $\Delta x$、時間刻み $\Delta t$ で離散化すると
\begin{equation}
  u_i^{n+1} = u_i^n + r\left(u_{i+1}^n - 2u_i^n + u_{i-1}^n\right),\qquad
  r = \frac{\kappa\,\Delta t}{\Delta x^2}
  \label{eq:scheme}
\end{equation}
となり、安定条件は $r \le 1/2$ である。

\section{方法}
\begin{enumerate}[label=(\arabic*)]
  \item 区間 $[0,1]$ を 50 分割し、初期条件 $u(x,0)=\sin\pi x$ を与える。
  \item $r=0.25,\ 0.5,\ 0.6$ の三通りで $t=0.1$ まで進める。
  \item 厳密解 $e^{-\pi^2\kappa t}\sin\pi x$ との最大誤差を記録する。
\end{enumerate}

\section{結果と考察}
$r=0.6$ では式~\eqref{eq:scheme} の解が振動して発散した。$r=0.25$ と $0.5$ では
誤差は $\Delta x^2$ に比例して減少し、理論どおりであった。

\section{結論}
陽的差分法は実装が容易だが、時間刻みの上限が厳しい。次回は陰的法（Crank--Nicolson）を試す。

\end{document}
